\subsection{Liberación de Recursos}

\subsubsection{Regla: la clase dueña de un recurso lo libera}

\textbf{Regla:} Todo objeto \texttt{IDisposable} creado y usado dentro de un mismo método se declara con \texttt{using}. Una clase que crea un objeto \texttt{IDisposable} y lo guarda en un campo es su dueña: implementa \texttt{IDisposable} y lo libera en su \texttt{Dispose}. Una dependencia recibida por constructor (Sección 10) no se libera, porque la clase no es su dueña (Microsoft, 2025g).

\textbf{Código correcto:}
\begin{lstlisting}[style=csharpstyle]
public sealed class ReplayRecorder : IDisposable
{
    private readonly FileStream _replayFile;

    public ReplayRecorder(string replayPath)
    {
        _replayFile = File.Create(replayPath);
    }

    public void Dispose()
    {
        _replayFile.Dispose();
    }
}
\end{lstlisting}

\textbf{Código incorrecto:}
\begin{lstlisting}[style=csharpstyle]
public sealed class ReplayRecorder
{
    private readonly FileStream _replayFile;

    public ReplayRecorder(string replayPath)
    {
        _replayFile = File.Create(replayPath);
    }
}
\end{lstlisting}

\subsubsection{Regla: cerrar un canal WCF con \texttt{Close} y \texttt{Abort}}

\textbf{Regla:} Un canal o cliente WCF no se libera con \texttt{using} ni con una llamada directa a \texttt{Close} dentro de \texttt{Dispose}. Se llama a \texttt{Close} dentro de un \texttt{try} y, si falla con \texttt{CommunicationException} o \texttt{TimeoutException}, se llama a \texttt{Abort}.

\textbf{Descripción:} \texttt{Close} puede lanzar una excepción cuando la red falla. Dentro del bloque \texttt{finally} que genera \texttt{using}, esa excepción oculta la que se lanzó en el cuerpo y omite el código que sigue; \texttt{Abort} libera el canal sin intentar comunicarse con el servidor (Microsoft, 2018).

\textbf{Código correcto:}
\begin{lstlisting}[style=csharpstyle]
public void Dispose()
{
    try
    {
        _channel.Close();
    }
    catch (CommunicationException)
    {
        _channel.Abort();
    }
    catch (TimeoutException)
    {
        _channel.Abort();
    }
}
\end{lstlisting}

\textbf{Código incorrecto:}
\begin{lstlisting}[style=csharpstyle]
public void Dispose()
{
    _channel.Close();
}
\end{lstlisting}
